% Overleaf-ready three-page research progress paper.
% Audience: company stakeholders, research supervisors, and technical partners.
% Compile with pdfLaTeX. All figures are drawn from verified project artifacts;
% no external citations are asserted in this progress brief.
\documentclass[10pt,a4paper]{article}

\usepackage[T1]{fontenc}
\usepackage[utf8]{inputenc}
\usepackage{lmodern}
\usepackage[margin=12mm]{geometry}
\usepackage{microtype}
\usepackage{booktabs}
\usepackage{tabularx}
\usepackage{array}
\usepackage{enumitem}
\usepackage{xcolor}
\usepackage{hyperref}
\usepackage{amsmath}
\usepackage{titlesec}

\definecolor{navy}{HTML}{17365D}
\definecolor{bluegray}{HTML}{EAF0F6}
\definecolor{greenbox}{HTML}{EAF4EA}
\definecolor{amberbox}{HTML}{FFF4D6}
\hypersetup{colorlinks=true,linkcolor=navy,urlcolor=navy}

\setlength{\parindent}{0pt}
\setlength{\parskip}{1.5pt}
\setlist{nosep,leftmargin=1.4em}
\setlength{\tabcolsep}{4pt}
\renewcommand{\arraystretch}{1.06}
\renewcommand{\baselinestretch}{0.96}
\titlespacing*{\section}{0pt}{6pt}{2pt}
\titlespacing*{\subsection}{0pt}{4pt}{1pt}

\newcommand{\status}[1]{\texttt{\detokenize{#1}}}
\newcommand{\infobox}[2]{%
  \begin{center}
  \colorbox{#1}{\begin{minipage}{0.95\textwidth}\small #2\end{minipage}}
  \end{center}}

\title{\vspace{-1.5em}\textbf{Industrial Energy Forecasting and Flexibility Research}\\
\large Current Progress, Verified Achievements, Limitations, and Decisions Ahead}
\author{Author Name \quad | \quad MSc Digital Engineering, Bauhaus-Universit\"at Weimar\\
\small Prepared for company and research feedback}
\date{November 2026}

\begin{document}
\maketitle
\vspace{-1.5em}
\small

\infobox{bluegray}{\textbf{Purpose of this brief.} This document summarizes
what has been built and tested across two public industrial-energy datasets,
what the results actually establish, what did not work, and which questions
now require industrial or academic feedback. It is a progress report, not a
claim of plant deployment or commercial savings.}

\section*{1. Research Goal and Current Position}

The project asks a practical question: \emph{when do better forecasts and more
realistic uncertainty scenarios improve constrained industrial-energy
decisions?} The work progressed from data-quality forensics to leakage-safe
forecasting, probabilistic scenarios, physically constrained optimization, and
economic decision evaluation. The central lesson is that forecast quality and
decision value are related but not equivalent.

The technical pipeline is now mature enough to support a bounded conference
paper and company discussion. It is classified as
\status{READY_WITH_LIMITATIONS} and \status{CATEGORY_C_EXPLORATORY}: software,
chronology, provenance, and numerical checks are strong, but the flexibility
asset and tariff remain counterfactual and no real factory schedule has been
validated.

\section*{2. Datasets and Why They Matter}

\begin{tabularx}{\textwidth}{p{0.16\textwidth}p{0.35\textwidth}X}
\toprule
Dataset & Physical target and period & Research role and key caution\\
\midrule
KIT SPARK & Positive facility electricity demand; 15-minute means; 2021--2024;
forecast at 12:00 Europe/Berlin on D-1 & Data-quality validation, deterministic
and probabilistic forecasting, scenarios, and the economic decision study.
Raw availability is incomplete and PV overlaps only 2024.\\
Honda R\&D Europe, Offenbach & Signed public-grid exchange; positive import,
negative export; 15-minute resolution; 2018--2023, with primary labels mainly
in 2020--2023; forecast at 13:00 D-1 & Provenance reconstruction, historical
baselines, LightGBM, foundation models, and scenario diagnostics. The target
already nets PV and CHP and is not validated gross demand.\\
\bottomrule
\end{tabularx}

The targets are physically different. Raw kW error and euro outcomes must not
be compared as though SPARK demand and Honda signed exchange were the same
problem. For SPARK, raw five-second measurements were aggregated only when at
least 75\% of an interval was available. At this threshold, mean absolute
aggregation error was 0.0060 kW, versus 0.0416 kW at approximately 70\%
coverage. For Honda, primary labels retain original observations and exclude
imputed or filled intervals.

\section*{3. What Has Been Completed}

\begin{tabularx}{\textwidth}{p{0.18\textwidth}X p{0.18\textwidth}}
\toprule
Workstream & Completed output & Status\\
\midrule
Data foundation & Missingness audit, aggregation validation, Honda provenance
masks, signed-target semantics, DST-safe physical timestamps & Completed and
frozen\\
Forecast protocol & Leakage-safe issuance cutoffs, chronological folds,
quality policies, target and feature contracts & Completed and frozen\\
Baselines & Historical analogue and calendar baselines, including frozen B5
benchmarks & Completed\\
Deterministic ML & Site-specific LightGBM with development-only selection and
untouched chronological tests & Completed\\
Foundation models & Zero-shot Chronos-2 and TimesFM 3.0 benchmarked without
fine-tuning & Completed\\
Uncertainty & Quantile calibration plus independent and coherent residual
scenario generation & Completed with calibration limitations\\
Optimization & Synthetic F1 work buffer, signed grid accounting,
non-anticipative stochastic scheduling, feasibility and settlement checks &
Validated framework\\
Decision value & Deterministic versus independent- and coherent-scenario
scheduling on common-feasible SPARK days & Completed exploratory study\\
\bottomrule
\end{tabularx}

\newpage
\section*{4. Main Verified Results}

\subsection*{4.1 Forecasting performance}

\begin{tabularx}{\textwidth}{p{0.18\textwidth}Xrr}
\toprule
Site & Method & MAE (kW) & RMSE (kW)\\
\midrule
SPARK 2024 & Historical baseline B3 & 2.7096 & 4.2319\\
 & Historical ensemble B5 & 2.5713 & 4.1518\\
 & LightGBM & \textbf{2.3520} & \textbf{3.4659}\\
\midrule
Honda 2023 & B5 & 76.480 & 101.338\\
 & LightGBM M1-CHP & \textbf{71.513} & \textbf{96.974}\\
 & TimesFM 3.0 zero-shot & 95.431 & 129.221\\
 & Chronos-2 zero-shot & 103.672 & 139.959\\
\bottomrule
\end{tabularx}

For Honda, LightGBM improves B5 MAE by 4.967 kW (6.495\%), with a paired daily
difference of -5.122 kW and 95\% interval [-6.833, -3.514]. However, export
recognition remains weak: export recall is 29.752\% for LightGBM and 30.809\%
for B5. The zero-shot foundation-model result is a bounded negative benchmark
for two checkpoints; it does not test fine-tuning or domain adaptation.

For SPARK, globally calibrated quantile LightGBM achieves WIS 1.4743, 80\%
coverage of 80.44\%, and 90\% coverage of 89.91\%. Holiday uncertainty remains
poorly represented: 80\% and 90\% coverage fall to 50.35\% and 68.29\%.

\subsection*{4.2 Scenario quality}

Whole-day residual scenarios (S2) improve dependence-sensitive diagnostics over
independent interval scenarios (S1). On rolling SPARK 2023, the variogram score
falls from 1.238 to 0.402 and ACF RMSE from 0.538 to 0.185. On Honda 2022, the
variogram score falls from 25.040 to 8.256 and ACF RMSE from 0.498 to 0.211.
Marginal calibration remains imperfect: S2 80\% coverage is 69.74\% for SPARK
and 75.21\% for Honda. Better dependence therefore does not mean that the full
scenario distribution is calibrated.

\subsection*{4.3 Decision-value result}

The decision experiment uses a frozen hypothetical F1 electrical-work buffer:
$\pm5$ kW power, $\pm10$ kWh energy state, 15-minute intervals, exact terminal
recovery, a zero-export cap, and no import cap. It is not a battery, identified
machine, or measured production process. The population contains 273 complete
SPARK days in 2024; 191 are feasible under all policies.

\begin{tabularx}{\textwidth}{Xrrr}
\toprule
Policy & Counterfactual cost reduction (euro/day) & Violation days &
Common-feasible reduction (euro/day)\\
\midrule
P2 deterministic LightGBM & 1.999617 & 82 & \textbf{2.326373}\\
P3 independent scenarios & 1.800671 & 16 & 2.182422\\
P4 coherent scenarios & 1.747536 & \textbf{11} & 2.135182\\
\bottomrule
\end{tabularx}

On the 191 common-feasible days, P2 is cheaper than P4 by 0.191191 euro/day
(95\% day-bootstrap interval [-0.271593, -0.131960]) and cheaper than P3 by
0.143951 euro/day. P3 versus P4 is unresolved: -0.047240 euro/day, 95\%
interval [-0.122619, 0.009287], Holm-adjusted $p=0.131$. P4 nevertheless cuts
modelled violation days from 82 under P2 and 16 under P3 to 11. A pre-specified
all-day feasibility correction does not reverse the cost ranking.

\infobox{greenbox}{\textbf{Strongest achievement.} The project connects
leakage-safe forecasting, scenario dependence, physical feasibility, and
settlement in one reproducible comparison. The result is scientifically useful
because it separates two effects: coherent scenarios improve modelled
feasibility, while deterministic scheduling retains lower mean counterfactual
cost.}

\section*{5. Engineering and Reproducibility Achievements}

\begin{itemize}
  \item Forecast issuance and feature availability are enforced independently;
  future load, PV, CHP, weather, and target-derived regimes are excluded.
  \item Physical days preserve 92, 96, and 100 intervals across DST transitions.
  \item Stochastic policies use one shared first-stage schedule; scenarios are
  not optimized independently as if future information were known.
  \item Import/export signs, energy units, terminal recovery, and feasibility
  are audited outside the solver.
  \item The replication configuration was frozen before outcome evaluation;
  two reruns produced byte-identical deterministic outputs.
  \item The complete suite passes 233/233 tests, including 20 replication tests.
  Original forecast/scenario artifacts were not regenerated or modified.
\end{itemize}

\newpage
\section*{6. Drawbacks and What the Results Do Not Establish}

\begin{tabularx}{\textwidth}{p{0.26\textwidth}X}
\toprule
Limitation & Consequence\\
\midrule
Hypothetical flexibility asset & Capacities do not describe verified SPARK or
Honda equipment; this is not plant validation.\\
Assumed price-publication timing & Scheduling at 12:55 is
\status{POST_PUBLICATION_TIMING_ASSUMED}; observation-level publication
timestamps are absent.\\
Counterfactual wholesale exposure & Monetary values are not customer bill
savings, trading profits, or validated revenue.\\
Previously exposed SPARK 2024 & The decision replication is exploratory rather
than a new independent confirmatory test.\\
Common-feasible filtering & The 191-day primary population may be less difficult
than all 273 days; the corrected all-day analysis remains secondary.\\
Only SPARK economics & Honda contributes forecast and scenario evidence, but no
Honda economic claim is supported without verified prices and timing.\\
Imperfect scenario calibration & S2 improves dependence but undercovers nominal
80\% and 90\% intervals.\\
Limited external validity & Two public sites and one primary asset/tariff setup
cannot establish general industrial flexibility value.\\
\bottomrule
\end{tabularx}

The research also produced important negative findings. More zero-shot
foundation-model benchmarking is unlikely to be high value without a new domain-
adaptation question. More calibration variants did not solve holiday uncertainty.
Most importantly, improved energy/variogram scores did not produce lower mean
decision cost. This is not a failure of the project: it identifies where added
model complexity does and does not create value.

\section*{7. Where We Stand Now}

\infobox{amberbox}{\textbf{Current position.} The forecasting and software
questions are sufficiently mature for thesis and conference writing. The main
remaining uncertainty is not which algorithm to add; it is whether the
counterfactual asset, tariff, and information timing represent a decision that
an industrial partner considers operationally meaningful.}

\textbf{What is established}
\begin{itemize}
  \item Site-specific LightGBM is a strong deterministic benchmark on both
  datasets, while Honda export recognition remains weak.
  \item Whole-day residual sampling represents temporal dependence better than
  independent interval sampling on the reported diagnostics.
  \item Under the frozen SPARK decision problem, coherence reduces modelled
  violations but does not improve mean cost.
\end{itemize}

\textbf{What remains unverified}
\begin{itemize}
  \item real process-level flexibility capacities, recovery rules, and operator
  acceptance;
  \item customer-specific tariffs, network charges, price-publication timestamps,
  and commercial settlement;
  \item an independent external holdout and a third site with different physical
  flexibility;
  \item whether coherent scenarios add value under stronger state coupling,
  peak charges, CVaR, or validated recourse.
\end{itemize}

\section*{8. Feedback Requested from the Company}

\begin{enumerate}
  \item Is the F1 work-buffer abstraction recognizable for any real process, and
  what power, energy, ramp, and recovery limits would be credible?
  \item Which operational objective matters most: energy cost, capacity peaks,
  export avoidance, production risk, or a combination?
  \item Which tariffs, network charges, market deadlines, and publication times
  should a realistic case use?
  \item What cost or operational meaning should be assigned to a modelled
  zero-export-cap violation?
  \item Is anonymized operational validation, a third dataset, or prospective
  shadow scheduling possible without exposing confidential production data?
\end{enumerate}

\section*{9. Recommended Next Step}

Do not add another forecasting model by default. First, validate one decision
case with an industrial partner: define the asset, tariff, information timeline,
and value of feasibility before running a new economic experiment. If such
parameters cannot be validated, the present work should be presented honestly
as a reproducible exploratory study whose main contribution is demonstrating
that \emph{improved temporal dependence is not equivalent to improved decision
value}.

\vfill
{\footnotesize\textbf{Claim boundary.} All monetary outcomes in this brief are
synthetic counterfactual wholesale-cost differences. They are not customer
savings, trading profit, battery value, or evidence of Honda/SPARK plant
performance. Dataset and code redistribution remains subject to the respective
repository and licensing conditions.}

\end{document}
